\honest{Tsuyoku vs. the Egalitarian Instinct}{Eliezer Yudkowsky}{2007}

Hunter-gatherer tribes are usually egalitarian, at least for men; the all-powerful chief belongs mostly to farming societies. Among most of them a hunter who brings in a spectacular kill ``will carefully downplay the accomplishment to avoid envy.'' \nb{Richard Lee's account of the !Kung bears this out: the hunter says he saw ``just a little tiny one.'' In Lee's account the others enforce the modesty, to keep a good hunter from pride.}

Anyone who aims to do their best will sooner or later aim above the average. If you are ashamed of doing better than others, ``the median will forever be your concrete wall.'' And do you dare say you mean to do better than those below average? How prideful of you! \nb{The opening example shows public modesty, not private shame; the post itself later allows that public modesty may be wise. It gives no case of anyone stopped at the median by shame.}

Pride in doing better than others may not be healthy, but ``Personally I've found it to be a useful motivator, despite my principles, and I'll take all the useful motivation I can get.'' Competition may be zero-sum, but so is Go. In any case, it cannot be healthy to be ashamed of doing better. ``And besides, life is not graded on a curve'': the will to do better has no point where it turns into the will to do worse, and a race with no finish line has no medals. Run as fast as you can, ``without worrying that you might pull ahead of other runners''; if you ignore that possibility, it may happen. \nb{The principles are probably those of ``Twelve Virtues of Rationality,'' which says ``spare no thought for whether others are doing worse,'' and is also the source of ``life is not graded on a curve.'' The post offers pride in beating others and indifference to rank as two reasons, and does not say how they fit together.}

Sooner or later you will have ``fewer sins to confess.'' Sooner or later you will set out to fix a flaw most people have not fixed. It may be wise to play down your success in public: ``People may forgive a touchdown, but not dancing in the end zone.'' It is easier to disclaim your worth in public, as long as everyone knows it is not true, and it can be fun to display your modesty when everyone knows how much you have to be modest about. But do not stop there. Whisper to yourself, ``Tsuyoku, tsuyoku!'', set a higher target, and keep running, ahead or behind.
